\cotefile{37}
\scanpage{3}
SGA : Manuscrits à frapper

\diagram{\begingroup\small
\(\begin{array}{llll}
\text{\deleted{SGA 1}} & \text{\deleted{Index notations et terminologique:}} & & \text{\deleted{\uline{10} p}}\\
\text{SGA \unsure{4|2}} & \text{\tikz[remember picture,baseline=(sgaiv.base)]\node[inner sep=0pt] (sgaiv) {Exp IV (fin)};} & \text{Verdier} & \text{\uline{\unsure{50|30}}}\\
& \text{Exp V} & \text{Verdier} & 80\smash{\raisebox{-1.5em}{\handwritten{\(\left.\rule{0pt}{3.6em}\right\}\)}}}\smash{\raisebox{-3em}{\handwritten{\(\left.\rule{0pt}{6.6em}\right]\)}}}\\
& \text{Exp VI (fin)} & \text{Verdier} & 50\\
& \text{Exp VI B} & \text{S\unsure{a|\deleted{q}a}int Donat} & \unsure{50|10}\\
& \text{\shortstack[l]{Tables des matières pour les trois\\ volumes, index}} & & 30\\
& \text{Appendice à \unsure{XVII|\deleted{I}XVII}} & \text{Saint Donat} & 10\\
\text{SGA 5} & \text{Exp IV} & ? & 100\ ?\\
& \text{\tikz[remember picture,baseline=(sgavii.base)]\node[inner sep=0pt] (sgavii) {Exp VII};} & \text{Jouanolou} & \text{\uline{50}}\\
& \text{\tikz[remember picture,baseline=(sgaxi.base)]\node[inner sep=0pt] (sgaxi) {Exp XI};} & \text{Bucur} & \text{\uline{70}}\\
& \text{Exp XII} & \text{Bucur} & 100\smash{\raisebox{-0.65em}{\handwritten{\(\left.\rule{0pt}{2em}\right]\)}}}\\
& \text{Exp XIV} & \text{Grothendieck} & 20\\
& \text{\shortstack[l]{Introduction, table des matières,\\ index}} & & 15\\
\text{\deleted{SGA 6}} & \text{\deleted{Exp XI}} & \text{\deleted{Ferrand}} & \text{\deleted{70 ?}}\\
& \text{\deleted{Table des matières, index}} & & \text{\deleted{10}}\\
\text{SGA 7} & \text{Exp I à \deleted{\unsure{II|III}}IV} & ? & \unsure{120|100}\ ?\\
& \text{Exp \deleted{\unsure{V|IV}}V} & \text{Mme Raynaud} & \text{\tikz[remember picture,baseline=(raynaudcount.base)]\node[inner sep=2pt] (raynaudcount) {30};}\\
& \text{\tikz[remember picture,baseline=(sgarim.base)]\node[inner sep=0pt] (sgarim) {Exp VI};} & \text{Rim} & \text{\tikz[remember picture,baseline=(rimcount.base)]\node[inner sep=2pt] (rimcount) {\uline{70}};}\\
& \text{\shortstack[l]{Introduction, table des matières,\\ index}} & & 15
\end{array}\)
\endgroup}
\margin{\handwritten{--- \unsure{tel|Cl} \unsure{Saarbach}}}
\margin{\handwritten{--- \uline{écrire Bucur}}}
\margin{\handwritten{\uline{écr. Deligne}}}
% Pencil strokes are anchored to their table rows; counts themselves remain typed.
\handwritten{\figure{\begin{tikzpicture}[remember picture,overlay]
\draw ([xshift=-.8cm,yshift=-.45cm]sgaiv.south west) -- ([yshift=-.08cm]sgaiv.south west) -- ([yshift=-.08cm]sgaiv.south east);
\draw ([xshift=-1.05cm,yshift=-.3cm]sgavii.west) -- (sgavii.west);
\draw ([xshift=-.95cm,yshift=-.4cm]sgaxi.west) -- (sgaxi.west);
\draw ([xshift=-1.4cm,yshift=-.55cm]sgarim.west) -- (sgarim.west);
\draw (raynaudcount.north west) -- (raynaudcount.north east) -- (raynaudcount.south east) -- (raynaudcount.south west) -- cycle;
\draw (rimcount.north west) -- (rimcount.north east) -- (rimcount.south east) -- (rimcount.south west) -- cycle;
\draw ([xshift=-1pt]rimcount.north west) -- ([xshift=1pt]rimcount.north east) -- ([xshift=1pt,yshift=-1pt]rimcount.south east) -- ([xshift=-1pt,yshift=-1pt]rimcount.south west) -- cycle;
\end{tikzpicture}}}

NB Les chiffres soulignés correspondent à des exposés dont le manuscrit est disponible dès maintenant; ceux suivis d'un signe ? à ceux dont il n'est pas certain qu'ils seront rédigés. En plus de la frappe d'exposés, tables de matières, index, il y a pour chacun des cinq séminaires \unsure{envisagés} un certain travail de mise\deleted{\unsure{r}} au point (\unsure{correction|corection} de fautes de frappe et autres erreurs, insertion de notes de bas de page etc), représentant un travail variable d'un séminaire à l'autre. Il faudrait sans doute fixer un forfait pour chacun d'eux, pour la rénumération de la secrétaire, après examen avec elle des corrections à faire ...

Il y a actuellement dans les 250 pages prêtes pour la frappe. J'espère que dans les deux ou trois mois qui \unsure{viennent}, il y aura aussi les env. \unsure{200|300} pages qui manquent encore pour SGA 4 (qui sont déjà partiellement rédigés ...), ce qui ferait dans les 450 à 500 pages à frapper pour les cinq à six mois qui viennent. Cependant je ne peux le garantir, cela dépendra de Mr. Verdier, qui est à Brandeis cette année. Je me mets en rapport avec lui ces jours-ci.
% archive sheet number: 2

\scanpage{4}
\begin{flushright}
Massy le 3.11.1970\\
\end{flushright}
Chère Madame Michon,\\

Je m'aperçois que dans les exemplaires SGA d'archives que j'ai emportés de chez vous, il manque les suivants:
\[
\text{SGA 4 XVII}\quad,\quad\text{SGA 4 VI (première partie)}
\]

D'autre part j'ai besoin de ces exemplaires pour faire faire l'édition photooffset en préparation. Pourriez vous me les retrouver ?

Bien cordialement\\
% archive sheet number: 3

\scanpage{5}
\begin{flushright}
Massy le 3.11.1970\\
\end{flushright}
Cher Deligne,\\

Pourrais-tu me dire si tu as terminé de regarder l'exposé de Rim SGA 7 VI, et si oui, me l'envoyer avec tes commentaires (sinon, me dire si tu as l'intention de le regarder et quand) ?

J'ai demandé à Mlle Altazin, qui a tapé ton exposé SGA 4 XVIII, de te l'envoyer avec le manuscrit, pour que tu le corriges. Pourrais-tu me dire si tu l'as reçu et si tu as l'intention de faire les corrections ? Quand elles seront faites, ou si tu ne veux pas les faire, envoyes moi le texte au net stp.

As-tu eu des nouvelles de Ferrand pour SGA 6 XI ?

Bien cordialement\\
% archive sheet number: 4

\scanpage{7}
\mspage{(ix)}
de ces questions que j'avais fait ailleurs. Ainsi, je me suis trouvé travailler pendant trois ans dans une institution alors qu'elle participait à mon insu à un mode de financement que je considère comme profondément immoral et extrêmement dangereux. Etant à présent seul à avoir cette opinion parmi mes collègues à l'IHES, ce qui a condamné à l'échec mes efforts pour obtenir la suppression des subventions militaires du budget de l'IHES, j'ai pris la décision qui s'imposait et quitte l'IHES le 30 Septembre prochain, et suspends également toute collaboration scientifique avec cette Institution, aussi longtemps qu'elle continuera à accepter de telles subventions. Monsieur L. Motchane, directeur de l'IHES, s'est déclaré d'accord pour qu'à partir du 1\textsuperscript{er} Octobre 1970, l'IHES ne diffuse plus les textes mathématiques dont je suis auteur, ou faisant partie du Séminaire de Géométrie Algébrique du Bois Marie. Comme il a été dit plus haut, la diffusion de ce Séminaire va être assurée par la maison Julius Springer, dans la série des Lecture Notes. Je suis heureux de remercier ici la maison Springer et Monsieur K. Peters pour l'aide efficace et courtoise qu'ils m'ont apportée pour rendre possible cette publication, en se chargeant en particulier de la frappe pour la photooffset des nouveaux exposés rajoutés aux anciens séminaires, et des exposés manquants des séminaires incomplets.

\begin{flushright}
Massy, Août 1970\\
\end{flushright}
% archive sheet number: 6

\scanpage{9}
\begin{flushright}
\handwritten{16 Juillet .}\\
\end{flushright}
Cher Grothendieck,\\

Voici la nouvelle rédaction de \handwritten{4}.\handwritten{10}.\handwritten{6}.

En relisant le n\textsuperscript{o} \handwritten{9} de S.G.A \handwritten{4} IV j'ai trouvé les broutilles suivantes :

p. \handwritten{149} l. \handwritten{16} lire \(E'\handwritten{\vee} E'' = E\)

p. \handwritten{158} l. \handwritten{5} lire \(s\unsure{\varepsilon|\in}\operatorname{Hom}(e_E,G)\)

p. \handwritten{158} compléter le diagramme \handwritten{9}.\handwritten{3}.\handwritten{6}.\handwritten{2}

p. \handwritten{159} l. \handwritten{15} lire ``le foncteur \(\handwritten{\sigma}\) composé de

Dans les n\textsuperscript{o} \handwritten{10} et suivant j'ai trouvé

p. \handwritten{2} l. -\handwritten{7} lire ``doù l'assertion ''

p. \handwritten{14} l. \handwritten{5} lire \(\handwritten{\mathscr{H}}\mathrm{om}_A\)

p.\handwritten{14} l. \handwritten{17} itou .

\margin{\handwritten{\uline{demander}\\\uline{\unsure{Messing}}}}
\handwritten{\figure{\begin{tikzpicture}\draw (0,0) -- (0,1);\end{tikzpicture}}}
Enfin , désagréable surprise , je me suis aperçu qu'il manquait le n\unsure{\deleted{?}\handwritten{\textsuperscript{o}}|\handwritten{\textsuperscript{o}}} \handwritten{14} dans les notes que tu m'as envoyés (Modules sur un topos défini par recollement ). Je me souviens avoir rédigé ce numéro et je pense bien te l'avoir donné . En tout cas je ne l'ai pas dans mes papiers . J'espère que ce numéro ne s'est pas perdu .

\begin{flushright}
Bien à toi\\
\handwritten{\unsure{J. L.|J. V.}}\\
\end{flushright}
% archive sheet number: 8

\scanpage{10}
\uline{Exercice} \handwritten{4}.\handwritten{10}.\handwritten{6} . Soient \(S\) un \(\underline{U}\)-site dont la topologie est moins fine que la canonique , \(\underline{\handwritten{\unsure{M|\mathscr{M}}}}\) une partie de \(\mathrm{Fl}(S)\) possédant les propriétés suivantes :

 a) Les morphismes de \(\underline{M}\) sont quarrables (I \handwritten{10}.\handwritten{3}) et \(\underline{\unsure{\mathbb{M}|\deleted{S}M}}\) est stable par changement de base.

 b) \(\underline{\mathbb{M}}\) contient les flèches identiques et est stable par composition .

 c) Une flèche \(u:X\to Y\) telle qu'il existe une famille couvrante \(Y_i\to Y\) , avec les \(X\times_Y Y_i\to Y_i\) dans \(\underline{M}\) , est elle-même dans \(\underline{M}\) .

 d) Pour tout objet \(X\) , toute famille couvrante de \(X\) est raffinée par une famille couvrante \((f_i:X_i\to X)\) où les \(f_i\) sont dans \(\underline{M}\).

Pour tout objet \(X\) , \deleted{\unsure{considérons le site}}\inserted{notons} \(S(X)\) (``petit site de \(X\)'') \inserted{le site} \deleted{dont la catégorie sous-jacente est}\inserted{constitué par} la sous-catégorie pleine de \(S/X\) formée des objets dont le morphisme structural est dans \(\underline{M}\) , munie de la topologie induite par celle de \(S\) (III \handwritten{3}.\handwritten{1}).

1) Montrer que les familles couvrantes \((f_i:X_i\to X)_{i\unsure{\varepsilon|\in}I}\) où les \(f_i\) sont dans \(\underline{M}\) forment une prétopologie sur \(S\) (II \handwritten{1}.\handwritten{3}) qui engendre la topologie de \(S\) .

2) Pour toute flèche \(u:X\to Y\) de \(S\) , montrer que le foncteur changement de base par \(u\) est un foncteur continu (III \handwritten{1}.\handwritten{1}) de \(S(Y)\) dans \(S(X)\) . On notera
\[
S(u)^{\handwritten{*}}:S(Y)^{\handwritten{\sim}}\to S(X)^{\handwritten{\sim}}
\]
le foncteur correspondant entre les catégorie\handwritten{s} de faisceaux (III \handwritten{1}.\handwritten{2},iv)).

3) Définir une équivalence entre le topos \(S^{\handwritten{\sim}}\) et la catégorie des systèmes
\[
(F_X)\ (X\varepsilon\operatorname{ob}S)\quad,\quad(\handwritten{\varphi_u})\ (u\varepsilon\mathrm{Fl}\ S)
\]
formés d'objets \(F_X\varepsilon\operatorname{ob}S(X)^{\handwritten{\sim}}\) , et pour toute flèche \(u:X\to Y\) dans \(S\) , d'un morphisme \(\handwritten{\varphi}_u:S(u)^{\handwritten{*}}(F_Y)\to F_X\) , ces systèmes
% archive sheet number: 9

\scanpage{11}
étant soumis à une condition de transitivité pour un composé \(v_{\unsure{\circ|o}}u\) de flèches de \(S\) , et à la condition que \(u\unsure{\varepsilon|\in}\underline{M}\) implique que \(\handwritten{\varphi}_u\) soit un isomorphisme .

\handwritten{3}) Définir des foncteurs ``restrictions '' et `` prolongement ''
\[
\operatorname{Res}_X:(S_{\handwritten{/}X})^{\handwritten{\sim}}\to S(X)^{\handwritten{\unsure{\vee|\sim}}}\ ,\quad
\operatorname{Prol}_X:S(X)^{\handwritten{\sim}}\to(S_{\handwritten{/}X})^{\handwritten{\sim}}\ ;
\]
montrer que \(\operatorname{Res}_X\) commute aux limites inductives et aux limites projectives et que \(\operatorname{Prol}_X\) est \uline{pleinement fidèle} , son image essentielle étant formée des faisceaux \(F\) tels que \(\handwritten{\varphi}_u\) soit un isomorp\deleted{g}hisme pour toute flèche \(u\) de \deleted{\unsure{du|\(M\)}} \(S_{\handwritten{/}X}\) .

\handwritten{4}) Définir un isomorphisme d'adjonction faisant de \(\operatorname{Res}_X\) l'adjoint à droite de \(\operatorname{Prol}_X\) . En conclure qu'il existe un morphisme de topos
\[
f:S(X)^{\handwritten{\sim}}\to(S_{\handwritten{/}X})^{\handwritten{\sim}}
\]
tel que \(f^{\handwritten{*}}=\operatorname{Res}_X\) et que le foncteur \(f^{\handwritten{*}}\) admet un adjoint à gauche \(f_!=\operatorname{Prol}_X\) . Montrer que \(\operatorname{Prol}_X\) transforme les faisceaux abéliens en faisceaux abéliens . Montrer que \(f:S(X)^{\handwritten{\sim}}\handwritten{\longrightarrow}(S_{\handwritten{/}X})^{\handwritten{\sim}}\) est un \uline{plongement} de topos (\handwritten{9}.\handwritten{1}.\handwritten{1}) et que par suite ``le petit topos '' \(\handwritten{S(X)^{\sim}}\) est \deleted{\unsure{un sous-topos (9.1}} équivalent à un sous-topos du ``gros topos '' \((S_{\handwritten{/}X})^{\handwritten{\sim}}=S^{\handwritten{\sim}}_{\handwritten{/}X}\) .

\handwritten{5}) Soit \(F\) un faisceau abélien flasque de \(S^{\handwritten{\sim}}_{\handwritten{/}X}\) (V \handwritten{4}) . M\textsuperscript{O}ntrer que \(\operatorname{Res}_X F\) est \(S(X)\)-acyclique\inserted{\handwritten{(V 4.1)}} . En déduire que pour tout faisceau abélien \(F\) de \(S^{\handwritten{\sim}}_{\handwritten{/}X}\) et tout entier \(q\) , le groupe \(H^q(X,F)\) (cohomologie calculée dans \(S^{\handwritten{\sim}}_{\handwritten{/}X}\) ) est canoniquement isomorphe au groupe \(H^q(X,\operatorname{Res}_X F)\) (cohomologie calculée dans \(S(X)^{\handwritten{\sim}}\) ) .

\handwritten{6}) Déduire de 5) que pour tout faisceau abélien \(G\) \deleted{sur} de \(S(X)^{\handwritten{\sim}}\) et tout entier \(q\) le groupe \(H^q(X,G)\) \deleted{(resp.} (cohomologie calculée dans \(S(X)^{\handwritten{\sim}}\) ) est canoniquement isomorphe aux groupes \(H^q(X,f_{\handwritten{*}}G)\) et \(H^q(X,\operatorname{Prol}_X G)\) (cohomologie calculée dans \(S^{\handwritten{\sim}}_{\handwritten{/}X}\) ) .
% archive sheet number: 10

\scanpage{12}
\uline{Exercice} \handwritten{4}.\handwritten{10}.\handwritten{7} . On reprend les hypothèses et les notations de l'exercice \handwritten{4}.\handwritten{10}.\handwritten{6} . On suppose de plus que \(\underline{M}\) possède la propriété suivante :

(Et) Pour tout couple de flèches composables \(u\) et \(v\) de \(\mathrm{Fl}\ S\) , les \deleted{ces} propriétés \(u\circ v\unsure{\varepsilon|\in}\underset{\handwritten{\deleted{\unsure{m|M}}}}{\underline{M}}\) et \(u\varepsilon\underline{M}\) impliquent \(v\handwritten{\unsure{\in|\varepsilon}}\underline{M}\) .

1) Montrer que pour toute flèche \(u:X\to Y\) , le foncteur de changement de base par \(u\) est un morphisme \deleted{du topos} \(S(u)\) du site \(S(X)\) dans le site \(S(Y)\) . En d'autres termes , montrer que \(S(u)^{\handwritten{*}}:S(Y)\to S(X)\) est exact à gauche .

2) Montrer que \(\operatorname{Prol}_X:S(X)^{\handwritten{\sim}}\to(S_{\handwritten{/}X})^{\handwritten{\sim}}\) est exact à gauche . En déduire qu'il existe un morphisme de topos
\[
g:S^{\handwritten{\sim}}_{\handwritten{\deleted{\unsure{u|m}}/}X}\to S(X)^{\handwritten{\sim}}
\]
tel que \(g^{\handwritten{*}}=\operatorname{Prol}_X\) et \(g_{\handwritten{*}}=\operatorname{Res}_X\) . Montrer que \(g\) est une rétraction du plongement \(f\) i.e qu'on a \(g\circ f\handwritten{\simeq}\operatorname{id}_{S(X)}\) et qu'on a un morphisme ``d'adjonction'' (entre morphismes de topos ) \(\operatorname{id}_{S^{\handwritten{\sim}}_{\handwritten{/}X}}\handwritten{\longrightarrow} f\circ g\) .

\handwritten{3}) Montrer que les conclusion de \handwritten{1}) et \handwritten{2}) ne sont plus valables lorsqu'on ne fait pas l'hypothèse (Et) . (\deleted{p}Prendre pour \(S\) la catégorie des schémas \unsure{munie|\deleted{,}munie} de la topologie f.p.p.f. (fidèlement plate et de présentation finie \handwritten{[} S.G.A. \handwritten{3} IV \handwritten{]} ) et pour \(\underline{M}\) les morphismes \deleted{fidèlement} plats et de \unsure{p\'{r}ésentation|présentation} finie ).

\handwritten{4}) Acheter une médaille en chocolat pour les rédacteurs .
% archive sheet number: 11

\scanpage{13}
\begin{flushright}
Massy le 23.6.1971\\
\end{flushright}
\unsure{Cher} Verdier,\\

Merci pour ta lettre du 23 Mai. C'est dommage que \illegible{1} \unsure{ne} m'aies pas \unsure{ènvoyé|ênvoyé|envoyé} ce malheureux exer 4.10.6, cela a\unsure{u|n|r}\illegible{1} \unsure{permis} d'envoyer enfin à l'imprimeur le fascicule 1 de \unsure{SGA} 4. La personne qui a frappé le texte termine mainte\illegible{1} \unsure{elle}\deleted{\illegible{1}} partira en vacances et Springer rend la machine \unsure{qu'elle} avait en location, faute d'autres manuscrits. \unsure{Cela} remet donc la publication même de ce \unsure{fascimle|fascicule} sine die.

Bien cordialement\\
% archive sheet number: 12

\scanpage{15}
\begin{flushright}
Massy le 18.5.1971\\
\end{flushright}
Cher Verdier,\\

Je ne trouve pas une minute pour m'occuper de la réédition de SGA 4. Saavedra m'a un peu aidé, \unsure{pour} faire un index pour le Chap IV , par.9 à 13. Mais on est bloqué faute d'avoir les deux appendices de Deligne (assez, de points, image inverse d'un plat), qu'il faudrait donner à la frappe. Les\deleted{\unsure{r}}as-tu, et si oui, peux-tu les envoyer pourlla frappe ? D'autre part, peux-tu envoyer une version corrigée de l'exercice 4.10.6 (médaille chocolat), j'ai \unsure{entièrement} oublié ce qui était canulé. C'est un exercice qui sert tout le temps (!), i\deleted{\unsure{n|o}}l ne devrait donc pas être question de le vider. Il doit y avoir quelque chose de vrai à mettre à la place du faux ...

Si tu en as marre de t'occuper de la mise au point ``définitive'' de SGA 4, je ne pourrais \unsure{t\deleted{e}'en|teen} blâmer. Je voudrais alors des instructions sur ce qu'il faut faire - je suis disposé à le laisser incomplet et dégueulasse, plûtôt que d'avoir à \deleted{une} y consacrer un jour ou deux aux dépens du boulot survie, infiniment plus important. Pourrais-tu demander à Artin ce qu'il faut faire avec les exposés après VI - j'ai des tas de notes d'errata, footnotes améliorations diverses - d'en faire des instructions compréhensibles pour une dactylo demanderait des jours de travail. Veut-il que je lui envoie mon paquet de notes ?
% archive sheet number: 14

\scanpage{16}
\begin{flushright}
Massy le 3.11.1970\\
\end{flushright}
Cher Verdier,\\

Ne m'étant guère occupé de Math depuis trois mois, je suis un peu perdu pour SGA 4. Si je me rappelle bien, tu as rédigé ou es en train de rédiger les parties suivantes, qui sont exactement ce qui manque pour que SGA 4 soit complet:
\diagram{\(\begin{array}{ll}
\text{\deleted{IV,}} & \text{\deleted{par.9 et suivants}}\\
\text{V} & \\
\text{VI} & \text{par.5 et suivants}
\end{array}\)}

Si je me rappelle bien, \unsure{\deleted{IV}VI|\(\mathbb{VI}\)} était en fait terminé d'être rédigé, il fallait seulement y apporter quelques modifications dont on avait discuté avant ton départ. De plus, je pense que tu as avec toi l'exposé
\[
\text{VI B}
\]
de Saint Donat, et je t'envoie également l'Appendice à XVII du même, dont j'avais apparemment lu les premières pages. P\deleted{\unsure{o}}ourrais tu me le renvoyer (ou le renvoyer à St Donat) avec tes annotations et commentaires ?

Tu dois avoir un\deleted{e} exemplaire du tirage de XVII; le XVIII n'a pas été tapé sur \(\mathbb{S}\)tencils, mais directement sur papier pour offset; la frappe est terminée, et Deligne est censé la corriger.

Ecris-moi stp où tu en es avec ta part de rédaction, et avec la lecture de VI B . \deleted{\unsure{Stinger|Springer}} Pour des raisons techniques, ce serait bien \unsure{agréable} si les textes V,VI,\(\mathrm{VI}_{\mathrm{B}}\) pouvaient être disponibles dans les deux, trois mois qui viennent. Dis-moi en tous cas si tu as l'intention de terminer, et si oui, quand tu penses avoir terminé.

Bien cordialement\\
% archive sheet number: 15
